\section{Fazit}
\label{sec:fazit}

Die Untersuchung zeigt, dass Sicherheitsprobleme im Scheduler nicht
nur durch offensichtlich unsicheren eigenen Code entstehen. Die reale
XSS-Schwachstelle lag im Zusammenspiel aus kontrollierbaren
Datenbankwerten und dem komfortablen, aber für diese Daten
ungeeigneten Standardrenderer von \texttt{vue-cal}. Der strukturelle
Fix über eigene Event-Slots und Vue-Interpolation beseitigt die
Ursache. Nach dem vollständigen Rollout auf \texttt{capstone-main}
und erneuter Prüfung besteht kein bekanntes aktuelles XSS-Risiko dieser Art.

Bei CORS war die Ausgangslage differenzierter: Die Spring-API war
nicht pauschal für fremde Origins freigegeben und die normale lokale
Anwendung vermeidet Cross-Origin-Aufrufe über den Vite-Proxy. Erst
die bewusst unsichere Demo-Konfiguration machte den öffentlichen
Datei-Export für eine fremde Browser-Origin lesbar. Die Demo
verdeutlicht damit korrekt die Wirkung von CORS, darf aber nicht mit
Authentifizierung oder einem automatischen Token-Leak gleichgesetzt werden.

Der größte verbleibende Handlungsbedarf liegt in den angrenzenden
Schutzschichten. Die globale \texttt{/files/**}-Freigabe muss
begrenzt, die Keycloak-Konfiguration gehärtet und der XSS-Fix durch
Regressionstests abgesichert werden. CORS sollte weiterhin
blockierend bleiben, solange kein fachlicher Cross-Origin-Bedarf
existiert. Die zentrale Erkenntnis aus beiden Themen ist gleich:
Sicherheitsgrenzen müssen explizit kontrolliert werden. Nutzerdaten
gehören an eine Textausgabegrenze; fremde Origins erhalten nur die
minimal erforderliche Lesefreigabe.
